We represent the polyhedron
\[
Q = \left\{
\begin{pmatrix}
x_1\\
x_2
\end{pmatrix}\in \mathbb{R}^2
| x_1 + x_2 \leq 1, x_1 - x_2 \leq 1, x_1 \geq 0, x_2 \geq 0 \right\}.
\]
in the following figure:
\begin{center}
\begin{tikzpicture}[scale=1.3]
  \begin{axis}[
  xlabel={$x_1$},
  ylabel={$x_2$},
  axis on top=true,
  axis equal,
  axis lines=middle,
  samples=41,
  xtick = {1},
  ytick = {-1,1},
  domain=0:1.5,
  ymin=-1.2,ymax=1.2,
  legend pos=south east,
]
\addplot [thick,color=gray!20,fill=gray!20, 
                    fill opacity=0.05]coordinates {
            (0, 1) 
            (1, 0)
            (0, 0)  };
\addplot[color=black, thin,domain=0:1.2] {1-x}
node[pos=0.5,sloped,above] {\tiny $x_1+x_2=1$};
\addplot[color=black,thin,domain=0:1.2] {x-1}
node[pos=0.5,sloped,below] {\tiny$x_1-x_2=1$};
\node [circle,fill=black,draw,scale=0.2,label=below left:$A$] at (0,0) {A};
\node [circle,fill=black,draw,scale=0.2,label=right:$B$] at (1,0) {B};
\node [circle,fill=black,draw,scale=0.2,label=below:$C$] at (0,1) {C};
\node [circle,fill=black,draw,scale=0.2,label=right:$D$] at (0,-1) {D};

\end{axis}
\end{tikzpicture}
\end{center}


The corresponding polyhedron in standard form is
$$\left\{
\begin{pmatrix}
x_1\\
x_2\\
x_3\\
x_4
\end{pmatrix} \in \R^4
| Ax = b, x \geq 0 \right\},$$
with
$$A =
\begin{pmatrix*}[r]
1 & 1 & 1 & 0\\
1 & -1 & 0 & 1
\end{pmatrix*},
b = 
\begin{pmatrix}
1 \\
1
\end{pmatrix},$$
with the variables $x_3$ and $x_4$ being slack variables.
Each basic solution is obtained by selecting 2 variables out of 4 to be in the basis. There is total of 6 possible selections of basic variables.\\
The matrix $B=(A_{j1},...,A_{jm})$ is composed of columns $j_1,...,j_m$ of the matrix $A$.


\begin{enumerate}

\item Basic solution with $x_1$ and $x_2$ in the basis ($j_1=1$,$j_2=2$).
\[
B = 
\begin{pmatrix*}[r]
1 & 1\\
1 &-1
\end{pmatrix*};
B^{-1} =
\begin{pmatrix*}[r]
\frac{1}{2} & \frac{1}{2}\\
\frac{1}{2} &-\frac{1}{2}
\end{pmatrix*};
x_B = B^{-1}b =
\begin{pmatrix}
1 \\
0
\end{pmatrix};
x =
\begin{pmatrix}
1 \\
0 \\
0 \\
0
\end{pmatrix}.
\]

This basic solution is feasible and corresponds to point $B=(1,0)$ in the figure.

\item Basic solution with $x_1$ and $x_3$ in the basis $(j_1=1,j_2=3)$.
\[
B = 
\begin{pmatrix*}[r]
1 & 1\\
1 &0
\end{pmatrix*};
B^{-1} =
\begin{pmatrix*}
0 & 1\\
1 & -1
\end{pmatrix*};
x_B = B^{-1}b =
\begin{pmatrix}
1 \\
0
\end{pmatrix};
x =
\begin{pmatrix}
1 \\
0 \\
0 \\
0
\end{pmatrix}.
\]

This basic solution is feasible and also corresponds to point $B=(1,0)$.

\item Basic solution with $x_1$ and $x_4$ in the basis $(j_1=1,j_2=4)$
\[
B = 
\begin{pmatrix}
1 & 0\\
1 &1
\end{pmatrix};
B^{-1} =
\begin{pmatrix*}
1 & 0\\
-1 & 1
\end{pmatrix*};
x_B = B^{-1}b =
\begin{pmatrix}
1 \\
0
\end{pmatrix};
x =
\begin{pmatrix}
1 \\
0 \\
0 \\
0
\end{pmatrix}.
\]

This basic solution is feasible and corresponds to point $B=(1,0)$ .

\item Basic solution with $x_2$ and $x_3$ in the basis $(j_1=2,j_2=3)$ .
\[
B =
\begin{pmatrix*}[r]
1 & 1\\
-1 &0
\end{pmatrix*};
B^{-1} =
\begin{pmatrix*}[r]
0 & -1\\
1 & 1
\end{pmatrix*};
x_B = B^{-1}b =
\begin{pmatrix*}[r]
-1 \\
2
\end{pmatrix*};
x =
\begin{pmatrix*}[r]
0 \\
-1 \\
2 \\
0
\end{pmatrix*}.
\]

This basic solution is not feasible because $B^{-1}b\ngeqslant0$. It corresponds to point $D=(0,-1)$ in the figure.

\item Basic solution with $x_2$ and $x_4$ in the basis $(j_1=2,j_2=4)$
\[
B = 
\begin{pmatrix*}[r]
1 & 0\\
-1 &1
\end{pmatrix*};
B^{-1} =
\begin{pmatrix}
1 & 0\\
1 & 1
\end{pmatrix};
x_B = B^{-1}b =
\begin{pmatrix}
1 \\
2
\end{pmatrix};
x =
\begin{pmatrix}
0 \\
1 \\
0 \\
2
\end{pmatrix}.
\]

This basic solution is feasible and corresponds to point $C=(0,1)$ in the figure.

\item Basic solution with $x_3$ and $x_4$ in the basis $(j_1=3,j_2=4)$ 
\[
B = B^{-1} =
\begin{pmatrix}
1 & 0\\
0 & 1
\end{pmatrix};
x_B = B^{-1}b =
\begin{pmatrix}
1 \\
1
\end{pmatrix};
x =
\begin{pmatrix}
0 \\
0 \\
1 \\
1
\end{pmatrix}.
\]

This basic solution is feasible and corresponds to point $A=(0,0)$ in the figure.
\end{enumerate}

In summary, out of the six basic solutions, three are degenerate, and
they correspond to the same vertex, identified by five active
constraints: two equality constraints, and three inequality
constraints: $x_2 \geq 0$, $x_3 \geq 0$, and $x_4 \geq 0$.
In the original polyhedron $Q$, among the six basic solutions,
\begin{itemize}
\item one corresponds to the vertex A, where the two
  constraints $x_1 \geq 0$ and $x_2 \geq 0$, associated with the non
  basic variables, are active,
\item three correspond to the vertex B, where the
  \textbf{three} constraints $x_2 \geq 0$, $x_1 + x_2 \geq 1$, and
  $x_1 - x_2 \geq
  1$ are active,
\item one corresponds to the vertex C, where the  two
  constraints $x_1 \geq 0$ and $x_1 + x_2  \leq 0$, associated with the non
  basic variables, are active,
\item one is infeasible, and corresponds to point D, where
  the two constraints $x_1\geq 0$ and $x_1 - x_2\leq 1$ are active.
\end{itemize}
The three basic solutions associated with vertex $B$ are
degenerate. From an algebraic point of view, one of the basic
variables is equal to zero. From a geometrical point of view, more
than two constraints are active. 

